% Περάστε στο πακέτο cseuoi-thesis τις κατάλληλες επιλογές για τη διατριβή σας
\documentclass[Greek,PhD]{cseuoi-thesis} % Διδακτορικό (στα Ελληνικά)
%\documentclass[English,MSc,Y]{cseuoi-thesis} % Μεταπτυχιακό με εξειδίκευση στα Υπολογιστικά Συστήματα (στα Αγγλικά)
%\documentclass[English,MSc,TH]{cseuoi-thesis} % Μεταπτυχιακό με εξειδίκευση στη Θεωρία Επιστήμης Υπολογιστών (στα Αγγλικά)
%\documentclass[English,MSc,L]{cseuoi-thesis} % Μεταπτυχιακό με εξειδίκευση στο Λογισμικό (στα Αγγλικά)
%\documentclass[English,MSc,E]{cseuoi-thesis} % Μεταπτυχιακό με εξειδίκευση στους Επιστημονικούς Υπολογισμούς (στα Αγγλικά)
%\documentclass[English,MSc,T]{cseuoi-thesis} % Μεταπτυχιακό με εξειδίκευση στις Τεχνολογίες - Εφαρμογές (στα Αγγλικά)


% Πακέτο για την εμφάνιση περιθωρίων (χρήσιμο για την εύρεση overfull boxes)
%\usepackage{showframe}

\usepackage{psfrag}
\usepackage{pst-all}
\usepackage{pst-node,pstricks}


% Ευρετήριο όρων
\makeindex[intoc]


% Συμπληρώστε τα στοιχεία σας στις παρακάτω εντολές (αφαιρώντας το \colorbox{gray}{})
\titleGr{\colorbox{gray}{Τίτλος Διατριβής}}
\titleEn{\colorbox{gray}{Thesis Title}}
\authorGr{\colorbox{gray}{Ιωάννης Παπαδόπουλος}}
\authorEn{\colorbox{gray}{Ioannis Papadopoulos}}
\Arthro{\colorbox{gray}{τον}}
\Aitiatiki{\colorbox{gray}{Ιωάννη Παπαδόπουλο}}
\dateGr{\colorbox{gray}{Φεβρουάριος 2016}}
\dateEn{\colorbox{gray}{February 2016}}
\advisorGr{\colorbox{gray}{Γεώργιος Νικολαΐδης, Επίκουρος Καθηγητής}}
\advisorEn{\colorbox{gray}{Georgios Nikolaidis, Assistant Professor}}


% Ρύθμιση του πακέτου hyperref
\makeatletter
\hypersetup{
	bookmarks=true,
	pdftitle={\@title},
	pdfauthor={\@author},
	pdfsubject={\csethesisType},
	colorlinks=true,
	linkcolor=black,
	citecolor=black,
	urlcolor=[HTML]{0000EE}
}
\makeatother


\begin{document}

% Σελίδες χωρίς αρίθμηση
\pagenumbering{gobble}

% Εκτύπωση της σελίδας τίτλου
\maketitle

% Αρχικοποίηση του minitoc
\dominitoc[n]

\chapter*{\cseafierwsi}
\label{ch:Dedication}

Η σελίδα αυτή είναι προαιρετική και περιέχει αφιέρωση σε κάποιο σημαντικό πρόσωπο.

\y\noindent Προτεινόμενο: 1-2 γραμμές.

\y\noindent Μέγιστο: 1 σελίδα.
 % Προαιρετικό
\include{FrontMatter/[Optional]Acknowledgements} % Προαιρετικό


% i, ii, iii, iv, ...
\pagenumbering{roman}

% Περιεχόμενα
\pdfbookmark[1]{\contentsname}{contents} % hyperref
\tableofcontents

% Εισαγωγή του ευρετήριου σχημάτων στα περιεχόμενα
\addstarredchapterc{\listfigurename} % minitoc

% Κατάλογος Σχημάτων
\listoffigures

% Εισαγωγή του ευρετήριου πινάκων στα περιεχόμενα
\addstarredchapterc{\listtablename} % minitoc

% Κατάλογος Πινάκων
\listoftables

% Εισαγωγή του ευρετήριου αλγορίθμων στα περιεχόμενα
\addstarredchapterc{\listalgorithmname} % minitoc

% Κατάλογος Αλγορίθμων
\listof{algorithm}{\listalgorithmname}

\chapter*{\csesymvolismoi}
\label{ch:Glossary}
% Εισαγωγή του κεφάλαιου στα περιεχόμενα
\addstarredchapter{\csesymvolismoi} %minitoc

Η σελίδα αυτή είναι προαιρετική.
Περιέχει ορισμούς και επεξηγήσεις εννοιών, όρων, συντομεύσεων, και συμβολισμών.
Αν είναι μεγαλύτερη από δύο σελίδες τότε πρέπει να πάει στο τέλος της διατριβής, αμέσως μετά τα παραρτήματα. % Προαιρετικό

\chapter*{Abstract}

Surname, Name, Initials. Title Abbreviation (MSc ή PhD), Computer
Science Department, University of Ioannina, Greece. Graduation
Month, Graduation Year. Tilte of Dissertation in English.  Thesis
Supervisor:  Full Name.




\y\noindent Maximum: 4 pages

\y\noindent Recommended: 2 pages

\chapter*{Εκτεταμένη Περίληψη στα Ελληνικά}

Όνομα Επώνυμο του Πατρώνυμο και της Μητρώνυμο.\\
Μ.Δ.Ε. στην Πληροφορική, Μήνας και Έτος Αποφοίτησης.\\
Τμήμα Μηχανικών Η/Υ και Πληροφορικής, Πανεπιστήμιο Ιωαννίνων.\\
Τίτλος Διατριβής.\\
Επιβλέπων: Όνομα Επώνυμο.\\
\bigskip

\y\noindent Μέγιστο μέγεθος: 2 σελίδες

\y\noindent Προτεινόμενο μέγεθος: 1/2 σελίδα


% 1, 2, 3, 4, ...
\pagenumbering{arabic}

% Εισαγωγή των κεφαλαίων
\chapter{Εισαγωγή}
\label{ch:Introduction}


\section{Στόχοι}
\label{sec:Objectives}
Εδώ αρχίζει η κανονική αρίθμηση των σελίδων, ενώ παραπάνω η αρίθμηση είναι i, ii, iii, όπως φαίνεται.
Το πρώτο κεφάλαιο αποτελεί μία εισαγωγή στο αντικείμενο της διατριβής και πρέπει να υπάρχει.
Στο περιεχόμενό του πρέπει να συμπεριλαμβάνονται και οι επί μέρους στόχοι της διατριβής:
\begin{itemize}
	\item Να αναλυθεί η συμπεριφορά...
	\item Η ανάπτυξη...
	\item Να κατηγοριοποιηθούν...
	\item Να επιβεβαιωθούν...
\end{itemize}


\section{Δομή της Διατριβής}
\label{sec:Structure}
Η διατριβή περιέχει ν κεφάλαια.
\chapter{Οδηγίες για τη Μορφή της Διατριβής}
\label{ch:Instructions}

Η μορφή αυτή καθιερώθηκε το 2005 και ενημερώθηκε το 2016. 
Από το 2019 ενημερώθηκε εκ νέου και προσαρμόστηκε ώστε να καλύπτει το νέο 
πρόγραμμα ``Μηχανική Δεδομένων και Υπολογιστικών συστημάτων''.
Στο κείμενο αυτό περιγράφεται η μορφή της διατριβής και παρέχονται σύντομες οδηγίες για τη χρήση του προτύπου Latex.

Προτείνεται η χρήση του Latex για την μορφοποίηση του κειμένου και παρέχεται πλήρες πρότυπο που παράγει αυτόματα κείμενο στη σωστή μορφή.
Βρίσκεται στο \url{https://github.com/vvdimako/cseuoi-thesis} σε δύο παραλλαγές: μία για το παλιό μεταπτυχιακό πρόγραμμα Πληροφορικής και μία για το νέο πρόγραμμα στη Μηχανική Δεδομένων και Υπολογιστικών Συστημάτων.
% Περισσότερες λεπτομέρειες, εδω. Πρέπει να βάλουμε ένα τελικό release για Πληροφορική με τις λεπτομέρειες μορφοποίησης και ένα πρώτο για το νέο πρόγραμμα.

Μπορούν βέβαια να χρησιμοποιηθούν άλλες εφαρμογές στοιχειοθεσίας κειμένου (Word, Open-Libre Office, κλπ), αλλά δεν διατίθεται πλήρες πρότυπο για αυτές, παρά μόνο ένα απλό βοήθημα το οποίο έχει σημαντικές ατέλειες, δεν ακολουθεί πλήρως τις απαιτήσεις μορφοποίησης και χρειάζεται πολλές, «χειροκίνητες» αλλαγές για να μπορέσει το τελικό κείμενο να έρθει στη σωστή μορφή.
Είναι ευθύνη του Συγγραφέα να κάνει τις απαιτούμενες αλλαγές ώστε το τελικό κείμενο να ακολουθεί το πρότυπο μορφοποίησης του Τμήματος. 


\section{Διαδικασία Υποβολής Τελικού Αντίτυπου Διατριβής}
\label{sec:Submission}
Για τη διαδικασία υποβολής, ανατρέξτε στον κανονισμό και σε πληροφορίες στην 
ιστοσελίδα του Τμήματος.
(Οι οδηγίες αφαιρέθηκαν από το κείμενο αυτό για να μην υπάρχει η πληροφορία 
σε πολλά σημεία γιατί με τις τροποποιήσεις που συμβαίνουν είναι πιθανό να 
υπάρξουν διαφορές μεταξύ τους.)
%Ο φοιτητής θα πρέπει να φέρει τη διατριβή στη μορφή που περιγράφεται, να περάσει την εξέταση, και να κάνει όλες τις απαιτούμενες από τους εξεταστές μετατροπές.
%Έπειτα, η διαδικασία που πρέπει να ακολουθήσει είναι η εξής:
%\begin{itemize}
%    \item Να πάρει την έγκριση του επιβλέποντος για την τελική μορφή της διατριβής (υπογραφή στη φόρμα Δ3).
%    \item Να υποβάλλει τη διατριβή στη Σ.Ε.Μ.Σ. για την τελική έγκριση της μορφής (όπως περιγράφεται εδώ) και να προσκομίσει τη φόρμα Δ3 την οποία θα υπογράψει ένα μέλος της επιτροπής.
%    \item Να παραδώσει τη διατριβή σε μορφή PDF στο γραφείο Γ17 μαζί με τη φόρμα Δ3.
%    \item Να παραλάβει από το γραφείο Γ17 αντίτυπα του εξωφύλλου, τα οποία θα βάλει μπροστά από τη σελίδα τίτλου της διατριβής. H τελευταία σελίδα θα πρέπει να είναι χοντρό λευκό γυαλιστερό χαρτί.
%    \item Να παραδώσει τα 3 αντίτυπα για Μεταπτυχιακή Εργασία Εξειδίκευσης ή 5 αντίτυπα για Διδακτορική Διατριβή στο γραφείο Γ17 και να πάρει την υπογεγραμμένη φόρμα Δ3, την οποία θα παραδώσει και στη Γραμματεία μαζί με τα υπόλοιπα έγγραφα για την αποφοίτηση.
%\end{itemize}



\section{Διαμόρφωση Κειμένου}
\label{sec:Text}

\subsection{Βασικές οδηγίες}
\label{subsec:Basic}


Η διατριβή πρέπει να είναι τυπωμένη σε μονή όψη, ενώ το κείμενο πρέπει να είναι εντός των εξής περιθωρίων (ίδια και στις μονές και στις ζυγές σελίδες):
\begin{itemize}
	\item Top: 2.5 cm.
	\item Bottom: 3 cm.
	\item Left: 2.5 cm.
	\item Right: 2.5 cm.
\end{itemize}

Σε όλη τη διατριβή, εκτός από τα Σχήματα και τους Πίνακες (συμπεριλαμβανομένων των λεζάντων όμως), πρέπει να χρησιμοποιείται γραμματοσειρά GFS Didot για το κανονικό κείμενο και Ubuntu Mono για εντολές (ως typewriter font), μεγέθους 12 στιγμών.

Οι παράγραφοι θα πρέπει να είναι πλήρως στοιχισμένες (justified) και με διάστημα μεταξύ γραμμών (διάστιχο, line spacing) στις 1.5 γραμμές. % \renewcommand{\baselinestretch}{1.4}
Η εσοχή της πρώτης γραμμής της παραγράφου είναι 20 στιγμές (ή 0.7 cm), αλλά η πρώτη παράγραφος ενός κεφαλαίου, ενότητας, κλπ, δεν έχει εσοχή.
Δεν υπάρχει επιπλέον κενό διάστημα πρίν ή μετά από κάθε παράγραφο (spacing before, after).
%ΑΕ: υπάρχει όμως διάστημα 20 στιγμών κάτω από κάθε ενότητα μέχρι την πρώτη παράγραφο ή μεταξύ τίτλων ενοτήτων όταν είναι συνεχόμενοι.

Οι όροι μπορούν να είναι σε \textit{πλάγια γράμματα}, 
χρησιμοποιώντας την εντολή \verb|\textit{πλάγια γράμματα}|, συνήθως την 
πρώτη φορά που χρησιμοποιούνται. Τα \textbf{έντο\-να γράμματα} πρέπει να 
χρησιμοποιούνται μόνο στην περίπτωση που είναι απαραίτητα για την κατανόηση 
του κειμένου, χρησιμοποιώντας την εντολή\\
\verb|\textbf{έντονα γράμματα}|.

Αν θέλετε να εισάγετε έναν νέο όρο στο ευρετήριο, θα πρέπει να χρησιμοποιήσετε 
την εντολή\index{νέος όρος} \verb|\index{νέος όρος}|. Για κάθε υποόρο που 
θέλετε να προστεθεί σε έναν προηγούμενο όρο, θα πρέπει να χρησιμοποιήσετε την 
εντολή\index{νέος όρος!υποόρος}\\
\verb|\index{νέος όρος!υποόρος}|.
% AE: Ποιά η διαφορά ευρετηρίου (μετά τα παραρτήματα) από το Γλωσσάρι (πριν την περίληψη);

Όταν θεωρείτε ότι είναι απολύτως απαραίτητο, μπορείτε να εισάγετε μία 
υποσημείωση με την εντολή \verb|\footnote{Αυτή είναι μία υποσημείωση.}|, 
η οποία εμφανίζεται στο κάτω μέρος της αντίστοιχης 
σελίδας\footnote{Αυτή είναι μία υποσημείωση.}.

To πακέτο \texttt{xgreek}\footnote{\url{https://www.ctan.org/pkg/xgreek?lang=en}.} 
ορίζει διάφορες χρήσιμες μακροεντολές, οι οποίες επιτρέπουν την εύκολη χρήση 
χαρακτήρων που είναι δύσκολα προσβάσιμοι από το πληκτρολόγιο.
Μερικές από αυτές τις μακροεντολές είναι οι εξής:
\begin{itemize}
	\item Άνω τόνος (\anwtonos): \verb|\anwtonos|.
	\item Άνω τελεία (\anoteleia): \verb|\anoteleia|.
	\item Σύμβολο του Ευρώ (\euro): \verb|\euro|.
	\item Σύμβολο τοις χιλίοις (\permill): \verb|\permill|.
\end{itemize}

Το κύριο αρχείο είναι το \texttt{SampleThesis.tex}, στο προοίμιο του οποίου 
θα πρέπει να περάσετε στο πακέτο \texttt{cseuoi-thesis} τις κατάλληλες 
επιλογές για τη διατριβή σας. Αν το κείμενο της διατριβής είναι στα Ελληνικά 
τότε θα πρέπει να περάσετε την επιλογή \texttt{gr}, ενώ αν το κείμενο της 
διατριβής είναι στα Αγγλικά τότε θα πρέπει να περάσετε την επιλογή \texttt{en}.
Για τη στοιχειοθεσία Μεταπτυχιακής Διπλωματικής Εργασίας θα πρέπει να περάσετε 
την επιλογή \texttt{msc}, ενώ για τη στοιχειοθεσία Διδακτορικής Διατριβής 
θα πρέπει να περάσετε την επιλογή \texttt{phd}.
Στην περίπτωση της Διδακτορικής Διατριβής δεν περνάτε κάποια άλλη επιλογή, ενώ
στην περίπτωση της Μεταπτυχιακής Διπλωματικής Εργασίας θα πρέπει επιπλέον να περάσετε 
την κατάλληλη επιλογή για μία από τις εξής ειδικεύσεις:
\begin{itemize}
	\item Προηγμένα Υπολογιστικά Συστήματα: \texttt{systems}.
	\item Επιστήμη και Μηχανική Δεδομένων: \texttt{data}.
\end{itemize}
% ΑΕ: Μετακίνησα την πρόταση πρίν το itemize γιατί έμοιαζε σαν νέα παράγραφο χωρίς indentation!

Στη συνέχεια του ίδιου αρχείου θα πρέπει να συμπληρώσετε τα στοιχεία σας 
στις αντίστοιχες εντολές, αφαιρώντας την εντολή \verb|\colorbox{gray}{}|.
Ειδικότερα, τα στοιχεία που θα πρέπει να συμπληρώσετε είναι ο τίτλος της 
διατριβής, το ονοματεπώνυμο του φοιτητή, ο μήνας και το έτος αποφοίτησης, 
καθώς και το ονοματεπώνυμο και τη βαθμίδα του επιβλέποντος καθηγητή.
Τα παραπάνω στοιχεία θα πρέπει να τα συμπληρώσετε και στα Ελληνικά και στα 
Αγγλικά προκειμένου να χρησιμοποιηθούν κατάλληλα σε διάφορα σημεία της 
διατριβής, όπως η σελίδα τίτλου και οι σελίδες με τις περιλήψεις.
Αν κάποια τμήματα της διατριβής σας είναι σκιασμένα, είτε δεν έχετε 
συμπληρώσει τα αντίστοιχα στοιχεία σας είτε δεν αφαιρέσατε την εντολή 
\verb|\colorbox{gray}{}| όταν τα συμπληρώσατε.

 Η σειρά των προκαταρκτικών τμημάτων: σελίδα τίτλου - εξώφυλλο,
εξεταστική επιτροπή, αφιέρωση, ευχαριστίες,% περιεχόμενα, κατάλογος σχημάτων, κατάλογος πινάκων, κατάλογος αλγορίθμων, γλωσσάρι, περίληψη, extended abstract,
κλπ. θα πρέπει να είναι όπως στο παρόν έγγραφο.
 Στη συνέχεια ακολουθεί το κυρίως κείμενο, οργανωμένο σε κεφάλαια και μετά τυχόν παραρτήματα,
ευρετήριο όρων, δημοσιεύεσεις συγγραφέα και σύντομο βιογραφικό σημείωμα του συγγραφέα.

\subsection{Αρίθμηση σελίδων}
\label{subsec:pageNumbering}

 Οι σελίδες πρίν τα «Περιεχόμενα» δεν αριθμούνται.
 Η αρίθμηση σελίδων από τα περιεχόμενα μέχρι και την σελίδα πρίν το πρώτο κεφάλαιο γίνεται με λατινική αρίθμηση (i, ii, κλπ).
 Στην συνέχεια η αρίθμηση ξεκινά από την αρχή και γίνεται με τη συνηθισμένη αραβική αρίθμηση (1, 2, κλπ), μέχρι και το τέλος του ευρετηρίου.
 Σημειώνεται ότι οι αριθμοί σελίδων, όπου υπάρχουν, βρίσκονται στο κάτω μέρος της σελίδας και στο κέντρο του κειμένου.


\subsection{Κεφάλαια, ενότητες κλπ.}
\label{subsec:Chapters}

Για να ξεκινήσετε ένα κεφάλαιο, χρησιμοποιείτε την εντολή 
\verb|\chapter{Τίτλος Κεφαλαίου}|. Παρόμοια, για μία ενότητα 
χρησιμοποιείτε την εντολή \verb|\section{Τίτλος Ενότητας}|, ενώ για υποενότητα 
την εντολή \verb|\subsection{Τίτλος Υποενότητας}|.
Είναι επιθυμητό να εισάγετε και μία ετικέτα κάθε φορά που χρησιμοποιείτε 
τις παραπάνω εντολές, το οποίο γίνεται με την εντολή \verb|\label{Ετικέτα}|, 
προκειμένου να μπορείτε να αναφέρεστε στο αντίστοιχο σημείο του κειμένου 
με την εντολή \verb|\ref{Ετικέτα}|.

Η μορφοποίηση του τίτλου του κεφαλαίου έχει ώς εξής:
Το κεφάλαιο ξεκινάει σε νέα σελίδα και από το επάνω μέρος της (επιπλέον του περιθωρίου) υπάρχει διάστημα 90 στιγμών και μετά η λέξη Κεφάλαιο και ο αριθμός κεφαλαίου, στοιχισμένα αριστερά, σε έντονα γράμματα τύπου small caps\footnote{Όλα γράμματα είναι κεφαλαία, αλλά τα ``πραγματικά κεφαλαία᾽᾽ π.χ. το 1ο γράμμα, είναι μεγαλύτερα.}, μεγέθους 18 στιγμών για τα μεγάλα και 14 για τα μικρά κεφαλαία γράμματα.
% ΑΕ στο .cls βλέπω ένα \vspace*{50\p@} αλλά στην πράξη βγαίνει πολύ μεγαλύτερο.
% ΑΕ επίσης το αληθινό μέγεθος των γραμμάτων μου μοιάζει 20pt !
Ακολουθεί διάστημα 35 στιγμών και ο τίτλος του κεφαλαίου στοιχισμένος δεξιά, σε έντονα γράμματα τύπου small caps, μεγέθους 18 και 14 στιγμών.
Από κάτω υπάρχει διάστημα 55 στιγμών και μετά ο μίνι πίνακας περιεχομένων του κεφαλαίου, που περιγράφεται παρακάτω.
% Το .cls λεει 40 πάλι.
Προσοχή, στα περιεχόμενα η λέξη «κεφάλαιο» δεν εμφανίζεται, υπάρχει μόνο ο αριθμός κεφαλαίου και ο τίτλος.


Οι τίτλοι κεφαλαίων, ενοτήτων, και οι λεζάντες σχημάτων κλπ, έχουν κεφαλαία γράμματα το πρώτο του τίτλου και μετά μόνο στα πρώτα γράμματα των Ουσιαστικών, Ρημάτων και Επιθέτων.
Αυτό πρέπει να γίνεται χειροκίνητα ακόμη και σε κείμενο γραμμένο σε Latex.

Ο μίνι πίνακας περιεχομένων ξεκινάει και τελειώνει με οριζόντιες γραμμές.
Αν δεν υπάρχουν ενότητες στο κεφάλαιο, τότε τοποθετείται μόνο μία οριζόντια γραμμή (για παράδειγμα βλ. Παράρτημα \ref{app:FirstAppendix}).
Στον πίνακα περιλαμβάνονται μόνο οι ενότητες πρώτου επιπέδου (π.χ. 2.1, 2.2, αλλά όχι 2.1.1) με την αριθμησή τους και με έντονους χαρακτήρες 12 στιγμών.
Δεν δίνονται οι αριθμοί σελίδας της κάθε ενότητας στον μίνι πίνακα.

Ο τίτλος των προκαταρκτικών τμημάτων (Αφιέρωση, κλπ), της Βιβλιογραφίας, Ευρετηρίου, κ.α. μορφοποιείται επίσης με έντονους χαρακτήρες, small caps 18 και 14 στιγμών,  αλλά με στοίχιση στα αριστερά και με μια οριζόντια γραμμή από κάτω.


Οι τίτλοι ενοτήτων πρώτου και δεύτερου επιπέδου αριθμούνται με τον αριθμό κεφαλαίου, αριθμό ενότητας, αριθμό υποενότητας, με τελείες ανάμεσα στους αριθμούς και μορφοποιούνται με έντονους χαρακτήρες 14 στιγμών.
Οι ενότητες πρώτου και δευτέρου επιπέδου εμφανίζονται στα περιεχόμενα με εσοχή ώστε ο αριθμός ενότητας - υποενότητας να είναι ευθυγραμισμένος με την αρχή του τίτλου της προηγούμενης ενότητας - κεφαλαίου (βλ. τα περιεχόμενα του παρόντος κειμένου).

Ενότητες τρίτου επιπέδου καλό είναι να αποφεύγονται. %ΑΕ: Απαγορεύονται?
Δεν έχουν αρίθμηση, δεν εμφανίζονται στα περιεχόμενα, και μορφοποιούνται με έντονους χαρακτήρες 12 στιγμών.


\subsection{Παραρτήματα}
\label{subsec:Appendices}
Προαιρετικά, μπορείτε να εισάγετε ένα ή περισσότερα παραρτήματα, τα οποία 
θα βρίσκονται μετά τη βιβλιογραφία και πριν το ευρετήριο.
Η μορφοποίησή τους είναι όμοια με τα κεφάλαια με τη διαφορά ότι η αρίθμηση των παραρτημάτων 
γίνεται με κεφαλαίους Ελληνικούς χαρακτήρες.


\section{Σχήματα}
\label{sec:Figures}
Τα σχήματα και οι λεζάντες τους πρέπει να είναι πάντα κεντραρισμένα και εντός 
των περιθωρίων του κειμένου. Για να εισάγουμε ένα σχήμα στο κείμενο, 
χρησιμοποιούμε τις παρακάτω εντολές:

\begin{verbatim}
\begin{figure}[t]
 \centering
 \includegraphics[width=0.65\textwidth]{Figures/ExponentialFunction.pdf}
 \caption{Η Εκθετική Συνάρτηση.}
 \label{fig:ExponentialFunction}
\end{figure}
\end{verbatim}

\begin{figure}[t]
	\centering
	\includegraphics[width=0.65\textwidth]{Figures/ExponentialFunction.pdf}
	\caption{Η Εκθετική Συνάρτηση.}
	\label{fig:ExponentialFunction}
\end{figure}

Όλα τα σχήματα, πίνακες, κλπ. πρέπει να είναι αριθμημένα και να αναφέρονται στον αντίστοιχο Κατάλογο Σχημάτων, Πινάκων, κλπ.
Η αρίθμησή τους ξεκινά από το 1 σε κάθε νέο κεφάλαιο.
Ο αριθμός σχήματος θα πρέπει να περιλαμβάνει τον αριθμό κεφαλαίου, και τον αριθμό σχήματος με μία τελεία μεταξύ τους.

Η αναφορά σε ένα σχήμα γίνεται με τον εξής τρόπο: 
«το αποτέλεσμα των παραπάνω εντολών φαίνεται στο 
Σχήμα~\ref{fig:ExponentialFunction}».
Προτείνεται η χρήση labels και cross-references αντί για αρίθμηση με το χέρι γιατί έτσι οι αλλαγές στη σειρά σχημάτων αλλάζουν αυτόματα και την αρίθμηση.
Στο Latex, χρησιμοποιούμε την εντολή 
\verb|Σχήμα~\ref{fig:ExponentialFunction}|.

Αν θέλουμε να εισάγουμε σε ένα σχήμα πολλές εικόνες μαζί, τότε χρησιμοποιούμε 
τις παρακάτω εντολές:

\begin{verbatim}
\begin{figure}[t]
 \centering
 \begin{subfigure}[t]{0.3\textwidth}
  \centering
  \includegraphics[height=0.15\textheight]{Figures/GraphA.pdf}
  \caption{}
  \label{subfig:GraphA}
 \end{subfigure}
 \hfill
 \begin{subfigure}[t]{0.3\textwidth}
  \centering
  \includegraphics[height=0.15\textheight]{Figures/GraphB.pdf}
  \caption{}
  \label{subfig:GraphB}
 \end{subfigure}
 \hfill
 \begin{subfigure}[t]{0.3\textwidth}
  \centering
  \includegraphics[height=0.15\textheight]{Figures/GraphC.pdf}
  \caption{}
  \label{subfig:GraphC}
 \end{subfigure}
 \caption{Τρία Γραφήματα.}
 \label{fig:ThreeGraphs}
\end{figure}
\end{verbatim}

\begin{figure}[t]
	\centering
	\begin{subfigure}[t]{0.3\textwidth}
		\centering
		\includegraphics[height=0.15\textheight]{Figures/GraphA.pdf}
		\caption{}
		\label{subfig:GraphA}
	\end{subfigure}
	\hfill
	\begin{subfigure}[t]{0.3\textwidth}
		\centering
		\includegraphics[height=0.15\textheight]{Figures/GraphB.pdf}
		\caption{}
		\label{subfig:GraphB}
	\end{subfigure}
	\hfill
	\begin{subfigure}[t]{0.3\textwidth}
		\centering
		\includegraphics[height=0.15\textheight]{Figures/GraphC.pdf}
		\caption{}
		\label{subfig:GraphC}
	\end{subfigure}
	\caption{Τρία Γραφήματα.}
	\label{fig:ThreeGraphs}
\end{figure}

Με τις παραπάνω εντολές εισάγαμε στο Σχήμα~\ref{fig:ThreeGraphs} τρία 
υποσχήματα, στα οποία μπορούμε να αναφερθούμε και ξεχωριστά αν θέλουμε, 
χρησιμοποιώντας τις αντίστοιχες ετικέτες που τους αναθέσαμε.


\section{Πίνακες}
\label{sec:Tables}
Με παρόμοιες εντολές μπορούμε να εισάγουμε και πίνακες. Για παράδειγμα, με τις 
παρακάτω εντολές δημιουργούμε τον Πίνακα~\ref{tab:Example}.

\begin{verbatim}
\begin{table}[t]
 \centering
 \caption{Ένας Πίνακας.}
 \label{tab:Example}
 \begin{tabular}{| l || l | l | l |}
  \hline
  κελί 1 & κελί 2 & κελί 3 & κελί 4\\
  \hline
  \hline
  κελί 5 & κελί 6 & κελί 7 & κελί 8\\
  \hline
  κελί 9 & κελί 10 & κελί 11 & κελί 12\\
  \hline
  κελί 13 & κελί 14 & κελί 15 & κελί 16\\
  \hline
  κελί 17 & κελί 18 & κελί 19 & κελί 20\\
  \hline
 \end{tabular}
\end{table}
\end{verbatim}

\begin{table}[t]
	\centering
	\caption{Ένας Πίνακας.}
	\label{tab:Example}
	\begin{tabular}{| l || l | l | l |}
		\hline
		κελί 1 & κελί 2 & κελί 3 & κελί 4\\
		\hline
		\hline
		κελί 5 & κελί 6 & κελί 7 & κελί 8\\
		\hline
		κελί 9 & κελί 10 & κελί 11 & κελί 12\\
		\hline
		κελί 13 & κελί 14 & κελί 15 & κελί 16\\
		\hline
		κελί 17 & κελί 18 & κελί 19 & κελί 20\\
		\hline
	\end{tabular}
\end{table}


\section{Αλγόριθμοι}
\label{sec:Algorithms}
Για τη στοιχειοθεσία αλγορίθμων σε μορφή ψευδοκώδικα, όπως φαίνεται στον 
Αλγόριθμο~\ref{alg:Example} για παράδειγμα, χρησιμοποιούμε τις παρακάτω εντολές:

\begin{verbatim}
\begin{algorithm}[t]
 \caption{Υπολογισμός $y = x^n$.}
 \label{alg:Example}
 \begin{algorithmic}[1]
  \REQUIRE $n \geq 0 \vee x \neq 0$
  \ENSURE $y = x^n$
  \STATE $y \leftarrow 1$
  \IF{$n < 0$}
  \STATE $X \leftarrow 1 / x$
  \STATE $N \leftarrow -n$
  \ELSE
  \STATE $X \leftarrow x$
  \STATE $N \leftarrow n$
  \ENDIF
  \WHILE{$N \neq 0$}
  \IF{$N$ is even}
  \STATE $X \leftarrow X \times X$
  \STATE $N \leftarrow N / 2$
  \ELSE[$N$ is odd]
  \STATE $y \leftarrow y \times X$
  \STATE $N \leftarrow N - 1$
  \ENDIF
  \ENDWHILE
 \end{algorithmic}
\end{algorithm}
\end{verbatim}

\begin{algorithm}[t]
	\caption{Υπολογισμός $y = x^n$.}
	\label{alg:Example}
	\begin{algorithmic}[1]
		\REQUIRE $n \geq 0 \vee x \neq 0$
		\ENSURE $y = x^n$
		\STATE $y \leftarrow 1$
		\IF{$n < 0$}
		\STATE $X \leftarrow 1 / x$
		\STATE $N \leftarrow -n$
		\ELSE
		\STATE $X \leftarrow x$
		\STATE $N \leftarrow n$
		\ENDIF
		\WHILE{$N \neq 0$}
		\IF{$N$ is even}
		\STATE $X \leftarrow X \times X$
		\STATE $N \leftarrow N / 2$
		\ELSE[$N$ is odd]
		\STATE $y \leftarrow y \times X$
		\STATE $N \leftarrow N - 1$
		\ENDIF
		\ENDWHILE
	\end{algorithmic}
\end{algorithm}


\section{Μαθηματικά}
\label{sec:Mathematics}

Οι μαθηματικές εκφράσεις απαιτούν ειδική μορφοποίηση.
Δεν επιτρέπεται η χρήση απλού κειμένου.
Για τη στοιχειοθεσία μαθηματικών εκφράσεων σε Latex χρησιμοποιούμε
τα παρακάτω περιβάλλοντα:
\begin{itemize}
	\item Εξίσωση: \verb|\begin{equation} ... \end{equation}|.

	\begin{equation}
		S_{n} = 1 + \sum_{k=1}^{n}\frac{1}{k^{2} + k}\,,\quad n \in \mathbb{N}\,.
		\label{eq:Example}
	\end{equation}

	\item Θεώρημα: \verb|\begin{theorem} ... \end{theorem}|.

	\begin{theorem}
		Τhe square of the hypotenuse (the side opposite the right angle) is equal 
		to the sum of the squares of the other two sides.
	\end{theorem}

	\item Λήμμα: \verb|\begin{lemma} ... \end{lemma}|.

	\begin{lemma}
		If a prime divides the product of two numbers, it must divide at least 
		one of those numbers.
	\end{lemma}

	\item Πόρισμα: \verb|\begin{corollary} ... \end{corollary}|.

	\begin{corollary}
		In any right triangle, the hypotenuse is greater than any one of the 
		other sides, but less than their sum.
	\end{corollary}

	\item Γεγονός: \verb|\begin{fact} ... \end{fact}|.

	\begin{fact}
		It takes 8 minutes 17 seconds for light to travel from the Sun’s surface 
		to the Earth.
	\end{fact}

	\item Σημείωση: \verb|\begin{remark} ... \end{remark}|.

	\begin{remark}
		This is a remark.
	\end{remark}

	\item Ορισμός: \verb|\begin{definition} ... \end{definition}|.

	\begin{definition}
		Addition is bringing two or more numbers (or things) together to make 
		a new total.
	\end{definition}

	\item Παρατήρηση: \verb|\begin{observation} ... \end{observation}|.

	\begin{observation}
		This is an observation.
	\end{observation}

	\item Aπόδειξη: \verb|\begin{proof} ... \end{proof}|.

	\begin{theorem}[Fermat's Last Theorem]
		There are no positive integers $x$, $y$, and $z$ that satisfy the 
		equation $x^{n} + y^{n} = z^{n}$ for any integer value of $n > 2$.
	\end{theorem}
	\begin{proof}
		``I have discovered a truly marvellous proof of this, which this margin 
		is too narrow to contain.''
	\end{proof}
\end{itemize}

Οι λέξεις Θεώρημα, Λήμμα, κλπ. και ο αριθμός τους είναι με έντονα γράμματα και βρίσκονται σε εσοχή 30 στιγμών από το αριστερό περιθώριο του κειμένου.
Το κείμενο του θεωρήματος, λήμματος, κλπ γράφεται με πλάγια γράμματα και ξεκινούν από την ίδια γραμμή, όχι από κάτω.
Η απόδειξη ξεκινάει με την λέξη «Απόδειξη.» σε πλάγια γράμματα, ακολουθεί το κείμενο της απόδειξης (ξεκινώντας από την ίδια γραμμή) και στο τέλος υπάρχει ένα άδειο τετραγωνάκι, στοιχισμένο δεξιά, όπως φαίνεται παραπάνω.

Τα παραπάνω είναι αριθμημένα με παρόμοιο τρόπο όπως τα σχήματα: καινούρια αρίθμηση ανά κεφάλαιο και ο αριθμός τους περιλαμβάνει αριθμό κεφαλαίου, τελεία και αύξοντα αριθμό αντικειμένου μέσα στο κεφάλαιο.
Μόνο οι εξισώσεις που δεν αναφέρονται πουθενά στο κείμενο, επιτρέπεται να μην έχουν αρίθμηση.
Οι υπόλοιπες εξισώσεις αριθμούνται, αλλά η αρίθμηση φαίνεται στα δεξιά της εξίσωσης και είναι μέσα σε παρενθέσεις.
Η αναφορά σε εξίσωση γίνεται ως εξής: στην Εξίσωση~\ref{eq:Example} φαίνεται $\ldots$
%ΑΕ: Δεν είμαι σίγουρος για τις παρενθέσεις στην αναφορά εξίσωσης. Δεν το κάνει αυτόματα όπως το έχουμε τώρα.

\section{Περιεχόμενα και Κατάλογοι}
\label{subsec:Contents}

  Τα περιεχόμενα περιλαμβάνουν όλους τους καταλόγους σχημάτων, πινάκων, αλγορίθμων, το γλωσσάρι, την περίληψη, το Extended Abstract, τα κεφάλαια με τα δύο πρώτα επίπεδα ενοτήτων τους, τη Βιβλιογραφία, τα παραρτήματα και το ευρετήριο.
  Ο αριθμός σελίδας όπου ξεκινάει το αντίστοιχο τμήμα βρίσκεται στοιχισμένο δεξιά.
  Για τις ενότητες μόνο, το κενό από το τέλος του τίτλου της ενότητας μέχρι τον αριθμό σελίδας καλύπτεται από τελείες, και δεν χρησιμοποιούνται έντονα γράμματα και αριθμοί.
  Αντίθετα, για τα υπόλοιπα τμήματα, χρησιμοποιούνται έντονα γράμματα-αριθμοί και δεν υπάρχουν τελείες μεταξύ του τίτλου και του αριθμού σελίδας.

  Οι καταχωρήσεις των καταλόγων σχημάτων, πινάκων και αλγορίθμων μορφοποιούνται όπως οι ενότητες του κειμένου.


\section{Διαχείριση Βιβλιογραφίας}
\label{sec:Bibliography}
Για τη δημιουργία της βιβλιογραφίας χρησιμοποιούμε το πακέτο BibTeX.
Για αυτό απαιτείται μία βιβλιογραφική βάση δεδομένων, η οποία αποθηκεύεται ως 
ένα απλό αρχείο κειμένου με κατάληξη bib.
Το αρχείο αυτό περιέχει καταχωρήσεις της παρακάτω μορφής:

\begin{verbatim}
@article{Newman2003a,
 author = {Newman, Mark E. J.},
 title = {The Structure and Function of Complex Networks},
 journal = {SIAM Review},
 volume = {45},
 number = {2},
 pages = {167--256},
 year = {2003},
 doi = {10.1137/S003614450342480}
}
\end{verbatim}

Κάθε καταχώρηση ξεκινά με τη δήλωση του τύπου της αναφοράς. Το παραπάνω 
παράδειγμα αποτελεί αναφορά σε ένα άρθρο περιοδικού, επομένως η καταχώρηση 
ξεκινά με τη δήλωση \verb|@article|. Στη συνέχεια αναθέτουμε ένα μοναδικό 
κλειδί στην καταχώρηση, π.χ. \verb|Newman2003a|, το οποίο χρησιμοποιούμε 
στο κείμενο της διατριβής για να αναφερθούμε σε αυτή με την εντολή 
\verb|\cite{Newman2003a}|. Τέλος, συμπληρώνουμε τα πεδία του αντίστοιχου 
τύπου αναφοράς, μερικά από τα οποία είναι υποχρεωτικά. Για παράδειγμα, στις 
καταχωρήσεις άρθρων είναι υποχρεωτική η συμπλήρωση των πεδίων \verb|author|, 
\verb|title|, \verb|journal|, και \verb|year|.

Η βιβλιογραφία της διατριβής στοιχειοθετείται αυτόματα μετά το τέλος των 
κεφαλαίων, με κάθε καταχώρηση να έχει έναν χαρακτηριστικό αριθμό.
Ο χαρακτηριστικός αριθμός της κάθε καταχώρησης εμφανίζεται μεταξύ αγκυλών 
στα σημεία του κειμένου της διατριβής όπου αναφερθήκαμε σε αυτή την καταχώρηση.
Για παράδειγμα, σε αυτήν την πρόταση αναφερόμαστε σε ένα άρθρο 
περιοδικού~\cite{Newman2003a}, σε μία εργασία συνεδρίου~\cite{DeCandia2007a}, 
σε μία τεχνική αναφορά~\cite{Jain1984a}, και σε ένα βιβλίο~\cite{Golumbic2004a}.

Το στυλ που ακολουθεί το πρότυπο είναι της IEEE Transactions.
Με το Latex η μορφοποίηση γίνεται αυτόματα.
Διαφορετικά θα χρειαστεί να ανατρέξετε στο
\url{https://ieee-dataport.org/sites/default/files/analysis/27/IEEE%20Citation%20Guidelines.pdf}.
για πληροφορίες για το πως μορφοποιείται κάθε είδος αναφοράς.





% Εισαγωγή της βιβλιογραφίας στα περιεχόμενα
\addstarredchapterc{\bibname} % minitoc

% Εισαγωγή της βιβλιογραφίας
\bibliographystyle{ieeetr}
\bibliography{Content/Bibliography}

% Προαιρετικά, μπορείτε να εισάγετε παραρτήματα
\appendix
\chapter{Τίτλος πρώτου παραρτήματος}
\label{app:FirstAppendix}
\minitoc

Εδώ είναι ο χώρος του πρώτου Παραρτήματος.

\begin{table}[h]
	\centering
	\caption{Πίνακας Παραρτήματος.}
	\label{tab:AppendixTable}
	\begin{tabular}{| l || l | l | l |}
		\hline
		κελί 1 & κελί 2 & κελί 3 & κελί 4\\
		\hline
		\hline
		κελί 5 & κελί 6 & κελί 7 & κελί 8\\
		\hline
		κελί 9 & κελί 10 & κελί 11 & κελί 12\\
		\hline
		κελί 13 & κελί 14 & κελί 15 & κελί 16\\
		\hline
		κελί 17 & κελί 18 & κελί 19 & κελί 20\\
		\hline
	\end{tabular}
\end{table}
\chapter{Τίτλος δεύτερου παραρτήματος}
\label{app:SecondAppendix}

\section{Τίτλος πρώτης ενότητας}
\label{sec:FirstSection}
Εδώ είναι ο χώρος της πρώτης ενότητας του δεύτερου Παραρτήματος.

\section{Τίτλος δεύτερης ενότητας}
\label{sec:SecondSection}
Εδώ είναι ο χώρος της δεύτερης ενότητας του δεύτερου Παραρτήματος.
\chapter{Τίτλος τρίτου παραρτήματος}
\label{app:ThirdAppendix}

Εδώ είναι ο χώρος του τρίτου Παραρτήματος.

\begin{figure}[h]
	\centering
	\includegraphics[width=0.125\textwidth]{Figures/bird.eps}
	\caption{Εικόνα Παραρτήματος.}
	\label{fig:AppendixFigure}
\end{figure}

% Εκτύπωση του ευρετηρίου
\printindex


% Σελίδες χωρίς αρίθμηση
\pagenumbering{gobble}

\chapter*{\csedimosieuseis}
% Εισαγωγή του κεφάλαιου στα περιεχόμενα
\addstarredchapter{\csedimosieuseis} % minitoc

Προαιρετικά, βάζουμε μία λίστα με τις δημοσιεύσεις του συγγραφέα. % Προαιρετικό

% Σύντομο Βιογραφικό
\include{BackMatter/ShortBio}

\end{document}
